\begin{textsample}{POS Dim 1 – 2009\_c6 – Score -2.00 – 2009\_c6/enade\_2009\_teatro\_7552cca521.txt}  \label{ex:f1_pos_228}
Capítulo 3 Análise Técnica da Prova Este capítulo tem por objetivo apresentar o desempenho dos estudantes de TEATRO no ENADE/2009.
Para isso, foram calculadas as estatísticas básicas da prova como um todo, bem como as estatísticas das questões relacionadas à Formação Geral e ao Componente Específico.
Nas Tabelas, são apresentadas as seguintes estatísticas : total da população, de presentes, média, erro-padrão da média, desvio-padrão, mediana, nota máxima, nota mínima e coeficiente de assimetria ( ver definição em 1.4.3 ).
Nas estatísticas das questões discursivas, foram detalhadas também as notas nulas : em branco, por protesto ou por insuficiência.
As estatísticas apresentadas neste capítulo contemplam, separadamente, os ingressantes, os concluintes e o total de estudantes, e foram calculadas tendo -se em vista as seguintes agregações : ( a ) Região Geográfica e País como um todo e ( b ) a Categoria Administrativa e a Organização Acadêmica.
As notas brutas das questões discursivas em cada uma das partes, Formação Geral/Componente Específico, foram calculadas a partir da média simples das notas de cada uma das questões correspondentes.
A primeira questão discursiva da Formação Geral foi desconsiderada e, portanto, a nota bruta referente às questões discursivas de Formação Geral foi igual à nota da questão mantida, a segunda.
Na componente específica, a nota bruta correspondeu à média das notas nas três questões.
As notas brutas das questões objetivas de cada uma das partes, Formação Geral/Componente Específico, foram calculadas a partir dos acertos dos estudantes considerando -se uma regra de três para a conversão para a escala de 0 a 100, descontando -se aquelas questões que foram desconsideradas pela comissão ou que não passaram pelo teste do índice de discriminação ponto-bisserial ( ver seção 1.4.6 para a descrição do índice ).
Questões com índices de discriminação inferior a 0,2 foram excluídas do cômputo da nota do estudante.
As notas brutas de Formação Geral e do Componente Específico para cada um dos estudantes foram calculadas como médias ponderadas das suas notas nas questões objetivas e discursivas, ponderação diferente para cada parte : C\_FG = C\_FG\textasciicircum{}OBJ × 0,6 + C\_FG\textasciicircum{}DIS × 0,4 C\_CE = C\_CE\textasciicircum{}OBJ × 0,85 + C\_CE\textasciicircum{}DIS × 0,15, onde C\_FG é a nota bruta da Formação Geral, C\_FG\textasciicircum{}OBJ é a nota bruta nas questões objetivas da Formação Geral, C\_FG\textasciicircum{}DIS é a nota bruta nas questões discursivas da Formação Geral, C\_CE é a nota bruta do Componente Específico, C\_CE\textasciicircum{}OBJ é a nota bruta nas questões objetivas do Componente Específico, e C\_CE\textasciicircum{}DIS é a A nota bruta da prova, C\_GER, por sua vez, é também uma média ponderada da nota de cada uma das partes, Formação Geral/Componente Específico : C\_GER = C\_FG × 0,25 + C\_CE × 0,75. 3.1 Estatísticas Básicas da Prova 3.1.1 Estatísticas Básicas Gerais A Tabela 3.1 apresenta as estatísticas básicas da prova por grupo de estudantes.
O total de estudantes inscritos para a prova foi de 2.838.
Destes, 72,9 \% compareceram, sendo que a presença foi maior entre concluintes ( 82,9 \% ) do que entre ingressantes ( 67,7 \% ).
A média geral da prova foi de 43,43, sendo que os ingressantes obtiveram média mais baixa, de 40,12, que os concluintes, cuja média foi 48,60.
O desvio-padrão geral foi de 16,9, sendo o do grupo de ingressantes menor ( 15,6 ), do que o do grupo dos concluintes ( 17,5 ), aqueles possuíram uma distribuição mais homogênea nas notas.
A nota máxima foi 89,1, obtida por pelo menos um concluinte ao passo que a maior nota obtida por um ingressante foi 86,1.
Os coeficientes de assimetria para os dois grupos considerados, o de concluintes e o de ingressantes, bem como para a população como um todo foram todos negativos, indicando uma maior concentração à direita da distribuição e um maior espalhamento à esquerda. nota bruta do Componente Específico, C\_CE\textasciicircum{}OBJ é a nota bruta nas questões objetivas do Componente Específico, e C\_CE\textasciicircum{}DIS é a nota bruta nas questões discursivas do Componente Específico.
A nota bruta da prova, C\_GER, por sua vez, é também uma média ponderada da nota de cada uma das partes, Formação Geral/Componente Específico : C\_GER = 0,25 × C\_FG + 0,75 × C\_CE. 3.1 Estatísticas Básicas da Prova 3.1.1 Estatísticas Básicas Gerais A Tabela 3.1 apresenta as estatísticas básicas da prova por grupo de estudantes.
O total de estudantes inscritos para a prova foi de 2.838.
Destes, 72,9 \% compareceram, sendo que a presença foi maior entre concluintes ( 82,9 \% ) do que entre ingressantes ( 67,7 \% ).
A média geral da prova foi de 43,43, sendo que os ingressantes obtiveram média mais baixa, de 40,12, que os concluintes, cuja média foi 48,60.
O desvio-padrão geral foi de 16,9, sendo o do grupo de ingressantes menor ( 15,6 ), do que o do grupo dos concluintes ( 17,5 ), aqueles possuíram uma distribuição mais homogênea nas notas.
A nota máxima foi 89,1, obtida por pelo menos um concluinte ao passo que a maior nota obtida por um ingressante foi 86,1.
Os coeficientes de assimetria para os dois grupos considerados, o de concluintes e o de ingressantes, bem como para a população como um todo foram todos negativos, indicando uma maior concentração à direita da distribuição e um maior espalhamento à esquerda.
Uma comparação entre o desempenho de ingressantes e concluintes pode ser realizada a partir dos dados apresentados no Gráfico 3.1.
Cumpre destacar que as notas dos concluintes, de um modo geral, foram mais elevadas do que as dos ingressantes, cuja distribuição de notas está à esquerda dos primeiros.
A ideia de melhor desempenho é reforçada pelas notas medianas dos dois grupos : 50,7 entre os concluintes e 40,5 entre os ingressantes.
Nas faixas de notas até 50 pontos o percentual de ingressantes é maior que o dos concluintes, situação esta que se inverte para as faixas de notas superiores a 50 pontos.
A exceção é o primeiro intervalo de 0 a 10 que inclui uma quantidade maior de concluintes do que de ingressantes.
Constatou -se que, aproximadamente 28 \% dos ingressantes e 51 \% dos concluintes situam -se nas faixas de notas superiores a 50,0 pontos.
O Gráfico 3.2 apresenta informações referentes ao desempenho geral de ingressantes e concluintes, comparando os resultados considerando a Região, Categoria Administrativa e Organização Acadêmica.
Levando -se em conta as notas médias dos estudantes em cada Região, observa -se que, nos dois grupos, há uma variação expressiva entre as Regiões.
As notas médias mais elevadas entre os concluintes ocorreram nas Regiões Sul ( 56,20 ) e Sudeste ( 48,96 ), ambas acima da média nacional ( 48,59 ).
Em relação aos ingressantes, as médias mais altas foram obtidas nas Regiões Sul ( 42,86 ) e Sudeste ( 41,14 ) nas demais Regiões a média foi inferior à média nacional ( 40,12 ).
Na Região Centro-Oeste foi obtida a média de notas mais baixa entre os ingressantes ( 36,51 ) e na Região Nordeste entre os concluintes ( 39,11 ).
Cursos nesta área são oferecidos em todas as esferas.
Levando -se em conta os agrupamentos dos estudantes em categorias administrativas, observa -se que a pontuação média mais elevada entre os concluintes e ingressantes foi encontrada em instituições de origem estadual ( respectivamente 54,55 e 42,60 ) – ambas acima da média nacional Tendo como foco as notas médias dos estudantes agrupados por Organização Acadêmica, notamos que centros universitários apresentaram as maiores médias para concluintes ( 52,13 ) enquanto IFECTS para ingressantes ( 45,73 ).
Ao passo que as piores notas médias foram observadas nas Faculdades tanto para ingressantes ( 40,71 ) como para concluintes ( 43,63 ). 3.1.2 Estatísticas Básicas de Formação Geral A Tabela 3.2 apresenta as estatísticas básicas de Formação Geral por grupo de estudantes.
O total de estudantes inscritos para a prova foi de 2838.
Destes, 72,9 \% do total compareceram, sendo que a presença foi maior entre concluintes ( 82,9 \% ) do que entre ingressantes ( 67,7 \% ).
A média geral da parte da prova de Formação Geral foi de 45,52, sendo que os ingressantes obtiveram média mais baixa, 44,39, que os concluintes, cuja média foi 47,29.
O desvio-padrão geral foi de 21,6, sendo o do grupo de ingressantes ( 21,4 ) menor do que o do grupo dos concluintes ( 21,8 ).
A nota máxima foi 100, obtida por estudantes de ambos os grupos.
Para esta parte, assim como para a nota geral, os coeficientes de assimetria são positivos para o total e para o grupo de ingressantes e negativo para o grupo de concluintes.
Como, porém, os valores são bem perto de zero, estes indicaram uma distribuição quase simétrica.
O Gráfico 3.3 permite a comparação entre concluintes e ingressantes relativamente ao desempenho na parte da prova que aborda a Formação Geral.
Nesta parte, se excluídas as provas com nota zero, tal como na prova como um todo, também os concluintes apresentaram desempenho superior, embora seja perceptível um maior equilíbrio.
A mediana dos concluintes foi igual à dos ingressantes, enquanto a média dos concluintes é maior do que a dos ingressantes.
Dos ingressantes, 40,3 \% alcançaram nota superior a 50 e dentre os concluintes, o percentual dos estudantes nesta faixa de notas é 47,3 \%.
No Gráfico 3.4 são apresentadas as informações referentes ao desempenho dos ingressantes e dos concluintes na parte de Formação Geral, em diferentes agregações como Região, Categoria Administrativa e Organização Acadêmica.
Considerando -se as notas médias dos estudantes segundo cada Região, observou -se que, tanto no grupo dos ingressantes quanto no grupo dos concluintes, houve diferenças marcantes.
As médias variam, entre os concluintes, de 38,63 na Região Nordeste a 52,77 na Região Sul.
Entre os ingressantes, os valores vão de 34,36 na Região Norte a 47,21 na Região Sudeste.
Nas Regiões Sul, Sudeste e Centro-Oeste, os concluintes apresentaram médias maiores do que os ingressantes, a situação inversa aconteceu no Nordeste.
A Região Norte não apresentou concluintes.
Tendo como foco as categorias administrativas, as maiores notas médias foram das instituições privadas para concluintes ( 49,57 ) e municipal para ingressantes ( 55,42 ), enquanto as das instituições municipais foram menores para concluintes ( 43,68 ) e estaduais para ingressantes ( 40,10 ).
Note que, com este recorte, ingressantes apresentaram médias de notas superiores aos dos concluintes nas instituições municipais.
Os estudantes das instituições privadas se apresentaram mais homogêneos do que os das demais instituições.
Considerando -se o tipo de Organização Acadêmica, Centros Universitários apresentaram os maiores valores para os concluintes ( 58,86 ) e IFECTs para os ingressantes ( 61,88 ). 3.1.3 Estatísticas Básicas do Componente Específico A Tabela 3.3 apresenta as estatísticas básicas referentes ao Componente Específico da prova, tendo sua média geral, de 42,78, inferior ao desempenho da prova de TEATRO como um todo, cuja média foi 43,43, como apresentado na Tabela 3.1, que por sua vez foi inferior ao desempenho na Formação Geral, que teve média igual a 45,52, conforme mostrado na Tabela 3.2.
Os concluintes apresentaram um desempenho médio ( 49,08 ) superior ao dos ingressantes ( 38,75 ), o oposto do ocorrido para a nota da Formação Geral, havendo, ainda, indícios de maior homogeneidade entre as notas médias dos ingressantes do que entre as dos concluintes, como sugeriu o desvio padrão, menor no grupo dos ingressantes do que no dos concluintes.
As notas máximas obtidas por concluintes e ingressantes foram iguais a 99,4 e 90,9, respectivamente.
As distribuições de notas para o total e o grupo dos concluintes se apresentaram com uma inclinação para notas maiores, como evidenciado pelos coeficientes de assimetria negativos.
Tendo ocorrido o oposto para o grupo dos ingressantes.
A situação de assimetria é confirmada pelo Gráfico 3.5 que proporciona uma comparação do desempenho entre concluintes e ingressantes em relação ao Componente Específico da prova.
Claramente a distribuição das notas médias dos ingressantes, à esquerda daquela dos concluintes, evidencia o percentual mais alto de notas médias nas faixas mais baixas para este grupo.
As medianas das distribuições das notas dos ingressantes e dos concluintes foram 38,1 e 50,2, respectivamente, enfatizando o melhor desempenho dos concluintes nesta parte da prova.
Assim como os Gráficos 3.2 e 3.4, o Gráfico 3.6 apresenta uma comparação dos resultados em relação às Regiões, Categoria Administrativa e Organização Acadêmica, agora levando em conta o desempenho dos ingressantes e dos concluintes no Componente Específico da prova.
Tendo como foco as notas médias dos estudantes, em determinada Região, observa -se que, entre os concluintes, as mais elevadas foram encontradas na Região Sul ( 57,40 ), valor acima da média nacional para esta parte ( 49,08 ).
Entre os ingressantes, as maiores notas ocorreram também nas Regiões : Norte ( 41,85 ), Sul ( 42,43 ) e Sudeste ( 39,17 ), todas acima da média nacional ( 38,75 ).
Novamente, para a análise segundo a Categoria Administrativa, todas as esferas estão representadas.
As médias mais elevadas, tanto entre os concluintes como entre os ingressantes, foram encontradas em instituições de origem estadual ( 56,53 entre os concluintes e 43,49 entre os ingressantes ) enquanto nas instituições privadas as médias foram ( 45,24 entre os concluintes e 35,33 entre os ingressantes ).
Quanto à Organização Acadêmica, Universidades, Centros universitários, Faculdades e Institutos Federais de Ciência e Tecnologia foram representados com este curso.
Considerando -se o tipo de Organização Acadêmica, para concluintes, as Universidades apresentaram os maiores valores, seguidas dos Centros Universitários e das Faculdades.
Já para o grupo dos ingressantes, os IFECTs apresentaram maiores valores que as Universidades. 3.2 Análise das Questões Objetivas 3.2.1 Formação Geral A Tabela 3.4 apresenta as estatísticas básicas relativas às oito questões objetivas da Formação Geral dos estudantes.
Adiantamos que a primeira questão foi desconsiderada pelo critério da correlação ponto-bisserial.
Como pode ser observado na Tabela 3.4, os estudantes concluintes obtiveram desempenho superior, com nota média de 61,69, em relação aos estudantes ingressantes cuja nota média foi 58,33.
A variabilidade entre os concluintes foi maior do que entre os ingressantes, como indicaram os desvios-padrão de 24,0, para o grupo de ingressantes e 24,8 para o grupo dos concluintes.
A amplitude total das notas foi a mesma para ambos os grupos : nota mínima de 0 e máxima de 100.
A mediana foi 57,1 para os ingressantes e 71,4 para os concluintes.
A assimetria, medida pelo coeficiente de assimetria ( skewness ) foi um pouco maior em módulo para os concluintes ( -0,7 ) do que para os ingressantes ( -0,3 ), ainda que ambos fossem negativos indicando uma assimetria para a esquerda.
A Tabela 3.5 apresenta a classificação das questões objetivas da Formação Geral, segundo o índice de facilidade ( questões em vermelho foram desconsideradas pelo critério ponto-bisserial ).
Como as questões foram comuns a todas as áreas, a tabela apresenta a classificação considerando -se os concluintes de todas as áreas e aqueles do curso de TEATRO.
De acordo com os índices obtidos, as questões objetivas da prova foram assim avaliadas para a área de TEATRO : das oito questões, três tiveram entre 60 \%, inclusive, e 85 \% de acertos e foram classificadas como fáceis ( somente uma para o total de concluintes ) ; outras duas foram consideradas medianas, alcançando entre 40 \%, inclusive, e 60 \% de acertos ( três para o total de concluintes ) ; duas obtiveram entre 15 \%, inclusive, e A Tabela 3.5 apresenta a classificação das questões objetivas da Formação Geral, segundo o índice de facilidade ( questões em vermelho foram desconsideradas pelo critério ponto-bisserial ).
Como as questões foram comuns a todas as áreas, a tabela apresenta a classificação considerando -se os concluintes de todas as áreas e aqueles do curso de TEATRO.
De acordo com os índices obtidos, as questões objetivas da prova foram assim avaliadas para a área de TEATRO : das oito questões, três tiveram entre 60 \%, inclusive, e 85 \% de acertos e foram classificadas como fáceis ( somente uma para o total de concluintes ) ; outras duas foram consideradas medianas, alcançando entre 40 \%, inclusive, e 60 \% de acertos ( três para o total de concluintes ) ; duas obtiveram entre 15 \%, inclusive, e 40 \% de acertos, sendo consideradas difíceis ( três para o total de concluintes ).
Por fim, apenas uma foi classificada como muito difícil obtendo percentual menor do que 15 \% de acertos ( uma para o total de concluintes ), coincidentemente a questão desconsiderada pelo critério ponto-bisserial.
Nenhuma questão foi classificada na categoria muito fácil.
A Tabela 3.6 descreve os resultados das análises das questões objetivas relativas à Formação Geral, segundo o poder de discriminação, utilizando -se para tal o índice de discriminação ponto-bisserial.
Como as questões foram comuns a todas as áreas, a tabela apresenta a classificação considerando -se os concluintes de todas as áreas e aqueles do curso de TEATRO.
A desconsideração da questão segundo este critério considerou o universo de todos os concluintes.
Duas das oito questões apresentaram índices iguais ou acima de 0,40 ( nenhuma para a população de concluintes como um todo ), sendo, portanto, classificadas como muito boas nesta característica, para este grupo de estudantes, uma foi classificada como boa com índice entre 0,30, inclusive, e 0,40 ( três para os concluintes como um todo ) e quatro foram classificadas como médias, com índice entre 0,20, inclusive, e 0,30 ( quatro para os concluintes como um todo ).
Uma das questões foi classificada como fraca tanto para a população de concluintes como um todo como para a de concluintes de TEATRO e foi desconsiderada pelo critério ponto-bisserial.
Na sequência, no Quadro 3.1, estão especificadas, por questão, as habilidades aferidas pelas questões de Formação Geral da prova de TEATRO, bem como os resultados da análise relativa aos índices de facilidade e de discriminação de cada questão.
O índice de facilidade para a área de TEATRO teve uma variação considerável, oscilando entre 0,09 até 0,85 e o de discriminação ficou entre 0,08 e 0,53.
A questão 8 que aferiu as habilidades de “ Ler e interpretar textos ”, “ Analisar e criticar informações ” ; ” Extrair conclusões por indução e/ou dedução ”, ” Estabelecer relações, comparações e contrastes em diferentes situações ”, “ Detectar contradições ”, “ Fazer escolhas valorativas avaliando consequências ”, “ Questionar a realidade ” e “ Argumentar coerentemente ” foi a que apresentou, para a população de concluintes de TEATRO, maior poder discriminatório com índice de 0,53.
A questão mais fácil, com 85 \% de acertos, foi também a de número 8.
De outra parte, a questão com menor índice de facilidade e de discriminação em Formação Geral foi a questão 1, que foi desconsiderada, e que deveria aferir as habilidades “ Ler e interpretar textos ”, “ Analisar e criticar informações ” ; ” Extrair conclusões por indução e/ou dedução ”, ” Estabelecer relações, comparações e contrastes em diferentes situações ”, “ Detectar contradições ”, “ Fazer escolhas valorativas avaliando consequências ”, “ Questionar a realidade ” e “ Argumentar coerentemente ”.
O Gráfico 3.7 mostra, a título de ilustração, o comportamento das respostas à questão 3 de Formação Geral para os estudantes concluintes da área de TEATRO como função da nota geral na prova.
Trata -se de uma questão considerada médio para os concluintes de TEATRO e médio para o conjunto de concluintes de todas as áreas.
O poder de discriminação foi médio para os estudantes da área e médio para o conjunto de todos os concluintes.
Neste Gráfico, cada uma das cinco curvas, representa o percentual de respostas em determinada alternativa da questão, em função do escore total dos estudantes concluintes da área de TEATRO.
A curva \textbf{roxa} corresponde à alternativa E, a correta para esta questão.
Assim, observa -se que entre os estudantes com notas mais baixas – desempenho geral mais fraco na prova como um todo –, houve a escolha de alternativas incorretas, especialmente a alternativa C – curva vermelha, seguida da alternativa D – curva em verde escuro, além da alternativa certa – curva vermelha.
A proporção de estudantes que selecionou a resposta correta, aumenta com o desempenho deste grupo de estudantes chegando ao acerto de 100 \% dos estudantes com notas pelo menos 87,5.
Esta análise permitiu verificar como a questão discriminou entre os grupos de desempenho.
Os Gráficos relativos às demais questões de Formação Geral constam do Anexo I. 3.2.2 Componente Específico A Tabela 3.7 apresenta as estatísticas básicas em relação às questões objetivas do Componente Específico.
Como pode ser observado, os estudantes concluintes obtiveram desempenho superior ( 49,88 ) ao dos estudantes ingressantes com uma média de 39,56 pontos.
Os desvios-padrão indicaram que a variabilidade entre os concluintes é maior ( 18,8 ) do que entre os ingressantes ( 16,9 ).
Os coeficientes de assimetria para os dois grupos foram negativos para os concluintes indicando uma concentração à esquerda e o inverso para os ingressantes.
A assimetria foi, porém, maior em modulo entre os concluintes.
A análise das questões objetivas do Componente Específico em TEATRO, em relação ao grau de facilidade da prova ( ver Tabela 3.8 ), mostrou que, das 27 questões, 8 encontraram -se na categoria difícil ( entre 15 \% e 40 \%, não incluso, de acertos ), 11 foram de nível médio ( entre 40 \% e 60 \%, não incluso, de acertos ), 2 estavam na categoria fácil ( entre 60 \% e 85 \%, não incluso, de acertos ) e nenhuma teve a classificação muito fácil ( a partir de 85 \% de acertos ).
Nenhuma questão teve percentual de acertos menor do que 15 \% ( classificação muito difícil ), 6 questões foram desconsideradas pela Comissão Assessora da Área de TEATRO.
Das 27 questões objetivas relativas ao Componente Específico ( ver Tabela 3.9 ), 6 obtiveram o índice de discriminação muito bom, 9 obtiveram o índice de discriminação bom e 4 obtiveram nível de discriminação médio.
O índice fraco que eliminaria a questão, foi alcançado por 2 questões e 6 questões foram desconsideradas pela Comissão Assessora da Área de TEATRO.
Na sequência, no Quadro 3.2 estão especificadas, por questão, as habilidades aferidas pelas questões do Componente Específico da prova, bem como os resultados da análise relativa aos índices de facilidade e de discriminação de cada questão.
O Gráfico 3.8 mostra, a título de ilustração, o comportamento da questão 12 do Componente Específico para os estudantes concluintes da área de TEATRO como função da nota geral na prova.
Trata -se de uma questão considerada médio para os concluintes de TEATRO.
O poder de discriminação foi bom para os estudantes da área.
Neste Gráfico, cada uma das cinco curvas representa o percentual de respostas em determinada alternativa da questão, em função da nota total dos estudantes concluintes de TEATRO.
A curva \textbf{roxa} correspondeu à alternativa E, a correta para esta questão.
Assim, observou -se que entre os estudantes com escores mais baixos – desempenho geral mais fraco na prova como um todo –, houve a escolha de alternativas incorretas, como a alternativa D – curva verde escuro ou a alternativa C – curva vermelha.
A proporção de estudantes que selecionou a resposta correta aumentou com o desempenho deste grupo de estudantes chegando ao acerto de 100 \% dos estudantes com notas pelo menos de 87,5.
Esta análise permitiu verificar como a questão discriminou entre os grupos de desempenho.
Os Gráficos relativos às demais questões do Componente Específico constam do Anexo I. 3.3 Análise das Questões Discursivas 3.3.1 Formação Geral A Tabela 3.10 apresenta as estatísticas básicas da prova por grupo de estudantes.
A análise dos resultados de desempenho dos estudantes nas questões discursivas que são apresentados na tabela 3.10 e no gráfico 3.9, mostra que as notas foram mais baixas no conjunto dessas questões que no das questões objetivas.
O total de estudantes inscritos para a prova foi de 2838.
Destes, 72,9 \% compareceram, sendo que a presença foi maior entre concluintes ( 82,9 \% ) do que entre ingressantes ( 67,7 \% ).
A média geral da prova foi de 24,44, sendo que os ingressantes obtiveram média mais baixa, de 23,60, que os concluintes, cuja média foi 25,77.
O desvio-padrão geral foi de 34,1 e tanto os concluintes quanto os ingressantes apresentaram um desvio-padrão similar 34,0 e 34,2, respectivamente.
A nota máxima foi 100,0, atingida por concluintes e ingressantes.
Os coeficientes de assimetria para os dois grupos considerados, o de concluintes e o de ingressantes, bem como para a população como um todo foram todos positivos, indicando uma maior concentração à esquerda da distribuição.
Para todos os grupos, cerca de 60 \% dos estudantes presentes obtiveram nota zero nas questões discursivas da Formação Geral.
Das 1275 provas com nota zero, 66 ( 5,2 \% ) continham um protesto, 317 ( 24,9 \% ) apresentaram a prova em branco e 892 ( 70,0 \% ) foram respondidas, porém de forma inadequada.
Entre os concluintes, a maioria ( 65,8 \% ) dentre aqueles com nota zero estavam de forma inadequada.
Entre os ingressantes, 72,4 \% das notas nulas ( zero ) decorreram do desempenho do próprio estudante verificadas na correção.
Uma comparação entre o desempenho de ingressantes e concluintes pode ser realizada a partir dos dados apresentados no Gráfico 3.9.
As notas dos concluintes, de um modo geral, foram ligeiramente mais elevadas do que as dos ingressantes, cuja distribuição de notas está à direita dos primeiros, reforçada pela proporção de notas nulas : 58,7 \% entre os concluintes e 63,5 \% entre os ingressantes.
A mediana foi igual a 0,0 para ambos os grupos.
Na sequência, discorrer -se-á, especificamente, sobre uma das questões discursivas de Formação Geral, a número 10 já que a número 9 foi desconsiderada.
Em especial, serão analisados os conteúdos englobados na questão e os resultados obtidos pelos estudantes.
Ao final, serão apresentados os comentários da banca corretora a respeito do observado nas respostas dos estudantes, suas impressões e conclusões.
Cumpre esclarecer que, tendo em vista que as questões discursivas de Formação Geral foram padronizadas, ou seja, constam de todas as provas, os comentários da banca foram os mesmos para todas as áreas acadêmicas, eis que direcionados a todos os estudantes que participaram do ENADE/2009. 3.3.1.1 Análise da Questão Discursiva 9 de Formação Geral DESCONSIDERADA 3.3.1.2 Análise da Questão Discursiva 10 de Formação Geral A Questão Discursiva nº 10 buscou avaliar as habilidades : • Ler e interpretar textos ; • Analisar e criticar informações ; • Extrair conclusões por indução e/ou dedução ; • Estabelecer relações, comparações e contrastes em diferentes situações ; • Fazer escolhas valorativas avaliando consequências ; • Questionar a realidade ; • Argumentar coerentemente.
O Quadro 3.3 apresenta as habilidades/competências, temas predominantes e estatísticas básicas da questão 10 de Formação Geral por grupo de estudantes.
Cumpre notar que como a questão 9 foi desconsiderada, a nota referente às questões discursivas da Formação Geral é idêntica à nota da questão 10.
A média geral da prova foi de 24,44, sendo que os ingressantes obtêm média mais baixa, de 23,60, que os concluintes, cuja média foi 25,77.
O desvio-padrão geral foi de 34,1 e tanto os concluintes quanto os ingressantes apresentam um desvio-padrão similar 34,0 e 34,2, respectivamente.
A nota máxima foi 100,0, atingida por concluintes e ingressantes.
Os coeficientes de assimetria para os dois grupos considerados, o de concluintes e o de ingressantes, bem como para a população como um todo foram todos positivos, indicando uma maior concentração à esquerda da distribuição.
Para todos os grupos, cerca de 60 \% dos estudantes presentes obtiveram nota zero nesta questão discursivas da Formação Geral.
Das 1275 provas com nota zero, 66 ( 5,2 \% ) continham um protesto, 317 ( 24,9 \% ) estavam em branco e 892 ( 70,0 \% ) foram respondidas, porém de forma inadequada.
Entre os concluintes, a maioria ( 65,8 \% ) dentre aqueles com nota zero estavam de forma inadequada.
Entre os ingressantes, 72,4 \% das notas nulas ( zero ) decorreram do desempenho do próprio estudante verificadas na correção.
Uma comparação entre o desempenho de ingressantes e concluintes para esta questão foi exatamente igual à já realizada para a nota discursiva de Formação Geral a partir dos dados apresentados no Gráfico 3.9. 3.3.1.3 Comentários sobre as respostas à Questão Discursiva 10 Apresentação A questão 10 de Formação Geral apresentou um questionamento capaz de analisar a habilidade do estudante em reconhecer o que é um direito social, optar por um direito social vigente que possa diminuir as desigualdades de renda familiar no Brasil e discutir os argumentos em favor do direito social escolhido como suporte à resposta.
Para análise das respostas apresentadas pelos estudantes, houve uma leitura inicial bastante criteriosa.
Foram considerados para efeito de pontuação, os seguintes argumentos : • Acesso à educação pública, gratuita e de qualidade, garantindo ao cidadão melhores condições de vida ; • Assistência à saúde, em seu contexto mais amplo, garantindo os cuidados essenciais de saúde, baseados em técnicas apropriadas ao alcance de todas as pessoas, \textbf{contando} com a participação da população ; • Permanência do estudante na escola, em todos os níveis escolares ( da educação infantil à educação superior, possibilitando a qualificação profissional ), o que garante melhores condições de trabalho e remuneração.
Considerou -se a qualificação e o desempenho docente, desenvolvendo nos estudantes melhores expectativas de aprendizagem e interesse pelo conteúdo da escola.
Os professores desenvolvendo metodologias adequadas e conhecimento pleno quanto a questões de aprendizagem ; • Condições dignas de trabalho, com remuneração que garanta qualidade de vida.
Foram consideradas as condições de trabalho, a preparação profissional para o trabalho e os cursos de aperfeiçoamento, oportunizando momentos de aprendizagem e desenvolvimento, a partir de situações desencadeadoras ; • Ser proprietário do imóvel em que reside, reduzindo os gastos com aluguel, garantindo melhores condições de vida, equilíbrio financeiro e construção do bem-estar familiar.
Desempenho dos estudantes na questão É importante salientar que havia respostas que desenvolveram adequadamente o questionamento apresentado, com fundamentação baseada em argumentos coerentes.
Contudo, em um número considerável de respostas, havia apenas um argumento e não dois, conforme fora solicitado.
Muitos estudantes informaram sobre problemas da educação, de um modo geral.
Algumas observações mencionadas pelos corretores durante o processo de correção foram : • há preocupação com uma educação de qualidade, com sugestões de cursos profissionalizantes e métodos preparatórios de aulas. • a educação é revelada como de baixa qualidade, exaltando -se treinamentos e cursos de atualização para a classe docente. • a saúde como essência para a vida.
O SUS é colocado como parte integrante de direito social, garantindo às pessoas fácil acesso. • oportunidade de emprego como forma de ascensão social, direito, cidadania. • os políticos brasileiros como causadores do despreparo educacional.
As respostas analisadas apresentaram norma culta, coesão e coerência razoáveis, que garantiram a compreensão da mensagem desenvolvida pelos estudantes.
Os erros ortográficos foram poucos e não comprometeram a discussão sobre o assunto e sobre os pontos de vista abordados.
Quanto à coerência, houve um desvio quanto ao objetivo da questão, já que alguns estudantes, através dos seus argumentos retrataram reivindicações, reclamações e situações de problemas enfrentados pelos brasileiros, citando golpes políticos.
Porém, mesmo quando se tratava desta avaliação, os estudantes foram autênticos através das explicações.
Houve erros quanto à concordância verbal e nominal.
Houve, também, dificuldade em utilizar os conectores corretamente.
Alguns estudantes, por falta de ideias, apresentaram em suas respostas redundância e prolixidade.
O desempenho foi considerado satisfatório, pois a maioria dos estudantes soube responder ao questionamento.
Houve domínio do conteúdo, pois as ideias, os argumentos e a originalidade corresponderam ao esperado. 3.3.2 Componente Específico Nas questões discursivas do Componente Específico – questões de números 38, 39 e 40 –, a diferença de desempenho entre os grupos não foi muito acentuada, inclusive menor do que para as questões discursivas de Formação Geral.
A Tabela 3.11 contém as estatísticas da média destas questões.
Observa -se que no grupo dos concluintes, a nota média alcançada foi 44,87, enquanto que no grupo dos ingressantes foi de 34,46, maior do que as obtidas nas questões discursivas de Formação Geral para os ingressantes ( 23,60 ) e para os concluintes ( 25,77 ) ( Tabela 3.10 ).
As médias foram também menores dos que as obtidas para as questões objetivas do Componente Específico ( Tabela 3.7 ) : respectivamente 39,56 e 49,88 para ingressantes e concluintes.
A análise de cada uma destas questões será feita nas subseções seguintes.
Os valores positivos para a medida de assimetria indicaram uma distribuição com concentração à direita.
Dos 2069 presentes, 46,9 \% obtiveram nota zero.
A tabela 3.12 apresenta um detalhamento da situação destes estudantes.
Mais da metade dos estudantes fizeram todas as questões.
Esta proporção foi ligeiramente maior entre os concluintes ( 51,4 \% ) do que entre os ingressantes ( 51,2 \% ).
A combinação de questões em branco e preenchidas correspondeu a 92,6 \% das provas.
Corresponderam a questões com protesto 3,9 \% no total das provas.
O Gráfico 3.10 apresenta a distribuição das notas nas questões discursivas do Componente Específico por grupo de estudantes.
A moda se situou no primeiro intervalo, aquele com valores entre 0 e 10 incluindo os extremos do intervalo, um máximo local é encontrado no último intervalo entre 90 e 100.
A impressão de assimetria informada pelos coeficientes de assimetria positivos ( Tabela 3.11 ) foi reforçada no gráfico pela concentração no primeiro intervalo. 3.3.2.1 Análise da Questão Discursiva 38 do Componente Específico A questão 38, cujas habilidades aferidas encontram -se descritas no Quadro 3.4, que ainda contém outras informações, seguido da Tabela 3.13 e do Gráfico 3.11, foi aquela em que os grupos de estudantes observados obtiveram um desempenho intermediário entre as questões discursivas do Componente Específico.
A nota média do grupo dos ingressantes foi de 36,42, e a dos concluintes 44,22.
O percentual de notas zero obtidas nos dois grupos foi de 56,5 \%, um pouco menor entre os concluintes ( 52,0 \% ) do que entre os ingressantes ( 59,3 \% ).
Mais de metade dos concluintes e dos ingressantes que tiraram zero responderam a questão ( ver Tabela 3.13 ).
Uma minoria ( 7,9 \% ) utilizou a questão para protestar.
Apesar do baixo desempenho, em ambos os grupos a nota máxima de 100,0 pontos foi alcançada.
O Gráfico 3.11 apresenta a distribuição das notas na questão 38 discursiva do Componente Específico por grupo de estudantes.
A moda se situou no primeiro intervalo, aquele com valores entre 0 e 10 incluindo os extremos do intervalo.
A impressão de assimetria informada pelos coeficientes de assimetria positivos ( Quadro 3.4 ) foi reforçada no gráfico pela concentração no primeiro intervalo.
A distribuição apresentou outro máximo local no ultimo intervalo, de notas mais altas, indicando uma certa polarização nas habilidades dos estudantes. 3.3.2.2 Comentários sobre as respostas à Questão Discursiva 38 As questões discursivas da Prova de Teatro ENADE/2009 de números 38, 39 e 40 foram de natureza semelhante.
As questões solicitavam a descrição de uma montagem hipotética sob um dos quatro pontos de vistas, a saber : encenação, atuação, cenografia e pedagogia do teatro.
A chave de correção proposta pela Equipe de Elaboração e Correção não foi aceita pela Comissão Assessora de Teatro do INEP/ Brasília ; portanto, foi sugerida uma outra forma de correção que desconsiderava os pontos de vistas acima mencionados em separado, massificando, assim, o olhar sobre as questões.
Qualquer que fosse o aspecto mencionado pelo estudante, considerando fatores como improvisação, jogo de cena, cenário, etc. deveria ser considerada como correta a questão em sua totalidade.
Dessa forma, o nível de dificuldade foi praticamente reduzido a nenhum.
Diante do exposto, os erros mais comuns consistiam em não mencionarem nada a respeito de uma montagem hipotética sobre os excertos dados.
Quanto ao aspecto da norma culta, coerência e coesão, houve uma incidência razoável de inadequações que foram pontuadas nas notas.
Merece destaque o fato de que um número considerável de estudantes optou pelo item “ Cenografia ”, \textbf{denotando} o palco italiano, não sendo este mencionado explicitamente.
Quanto ao item “ Pedagogia do Teatro ” houve despreparo por parte dos estudantes em relação a este aspecto.
Quanto ao item “ Atuação ” não houve referência aos jogos teatrais, linguagem de atuação.
O item “ Encenação ” \textbf{denota} que os estudantes abordaram a posição de atores no espaço, focos de luz, tipo de figurino, maneira de falar e o movimento corporal dos atores.
Alguns estudantes não apresentaram argumentos sobre as relações dialógicas.
CONCLUSÃO : Muitos estudantes encontram -se bem sucedidos como intérpretes e executantes cênicos, porém apresentam dificuldades ao assumir papéis de condutores e pesquisadores. 3.3.2.3 Análise da Questão Discursiva 39 do Componente Específico A questão 39, cujas habilidades a serem aferidas encontram -se descritas no Quadro 3.5, que ainda contém outras informações, seguido da Tabela 3.14 e do Gráfico 3.12, foi aquela em que os grupos de estudantes observados obtiveram o pior desempenho entre as questões discursivas do Componente Específico.
A nota média do grupo dos ingressantes foi de 30,70, e a dos concluintes 41,67, bem melhor entre os concluintes.
Isto ocorre, principalmente pelo grande percentual de notas nulas : 62,9 \% entre os ingressantes e 53,0 \% entre os concluintes.
Das provas com notas nulas nesta questão, a proporção deixada em branco foi de : 32,4 \% entre os ingressantes e 25,2 \% entre os concluintes ( ver Tabela 3.14 ).
Entre os com nota nula, uma pequena minoria ( 7,2 \% ) utilizou a questão para protestar.
Apesar do baixo desempenho, em ambos os grupos a nota máxima de 100,0 pontos foi alcançada.
Os coeficientes de assimetria positivos indicaram uma concentração à esquerda e espalhamento à direita.
O Gráfico 3.12 apresenta a distribuição das notas na questão 39 discursiva do Componente Específico por grupo de estudantes.
A moda se situou no primeiro intervalo, aquele com valores entre 0 e 10 incluindo os extremos do intervalo.
A impressão de assimetria informada pelos coeficientes de assimetria positivos ( Quadro 3.5 ) foi reforçada no gráfico pela concentração no primeiro intervalo.
A distribuição apresentou um máximo local no décimo intervalo, de notas mais altas, indicando uma certa polarização nas habilidades dos estudantes. 3.3.2.4 Comentários sobre as respostas à Questão Discursiva 39 Vide comentários da questão 38. 3.3.2.5 Análise da Questão Discursiva 40 do Componente Específico A questão 40, cujas habilidades aferidas encontram -se descritas no Quadro 3.6, que ainda contém outras informações, seguido da Tabela 3.15 e do Gráfico 3.13.
A questão foi aquela em que os grupos de estudantes observados obtiveram o melhor desempenho entre as questões discursivas do Componente Específico.
A nota média do grupo dos ingressantes foi de 36,31, e a dos concluintes 48,79.
Nesta questão ocorreu também o menor percentual de notas nulas : 57,2 \% entre os ingressantes e 45,6 \% entre os concluintes.
A maioria das provas com nota zero decorreu de provas preenchidas : 46,4 \% entre os ingressantes e 46,7 \% entre os concluintes ( ver Tabela 3.15 ).
Uma minoria ( 8,9 \% ), dentre aqueles com nota nula, utilizou a questão para protestar.
Em ambos os grupos a nota máxima de 100,0 pontos foi alcançada.
Os coeficientes de assimetria positivos entre os ingressantes indicaram uma concentração à esquerda e o coeficiente zero entre os concluintes, simetria na distribuição.
O Gráfico 3.13 apresenta a distribuição das notas na questão 40 discursiva do Componente Específico por grupo de estudantes.
A moda se situou no primeiro intervalo, aquele com valores entre 0 e 10 incluindo os extremos do intervalo.
A impressão de assimetria informada pelo coeficiente de assimetria positivo entre os ingressantes ( Quadro 3.6 ) foi reforçada no gráfico pela concentração no primeiro intervalo.
A distribuição apresentou um máximo local no décimo intervalo, aquele com notas entre ( 90,100 ], indicando uma certa simetria entre as notas não nulas e uma polarização das habilidades dos estudantes. 3.3.2.6 Comentários sobre as respostas à Questão Discursiva 40 Vide comentários da questão 38.

% matched lemmas: contar, denotar, roxa
\end{textsample}
